% Part: proof-theory
% Chapter: normalization
% Section: translations

\documentclass[../../../include/open-logic-section]{subfiles}

\begin{document}

\olfileid{pt}{nor}{tr}

\olsection{د عادي \Log{N2i}-\usetoken{P}{proof} او \Log{G2i} تر منځ اړول}

په \Log{N2i} کښې !!^{proof}s په \Log{G2i} کښې
!!{proof}s ته اړولے شو او برعکس هم۔ لومړے، د بې‌پرېکونه
\Log{G2i} !!{proof}s اړول عادي \Log{N2i}-!!{proof}s ورکوي۔
\textbf{د سرچينه‌يي نوم سمون OLCMP-177:} دلته کاپي شوے
نيمګړے نظام نوم د همدغه حدسي سيکوېنټ نظام په بشپړ نوم بدل شو۔

د دې نتيجې د بيانولو دپاره لږ اصطلاحي تيارے پکار دے، څو
په~\Log{N2i} او~\Log{G2i} کښې د سيکوېنټونو کيڼ سياقونه په
اسانۍ سره ونښلوو۔ د نښه لرونکو !!{formula}s سېټ~$\Gamma'$
له بې‌نښو !!{formula}s د څوګوني سېټ~$\Gamma$ سره
\emph{مطابقت لري} که د هر !!{formula}~$!A$ دپاره په~$\Gamma'$
کښې د $x : !A$ په بڼه د !!{element}s شمېر په~$\Gamma$ کښې
د~$!A$ له تکراري شمېر څخه کم يا ورسره برابر وي۔

\begin{cor}
که بې‌پرېکونه $\Log{G2i} \Proves \Gamma \Sequent \Delta$
وي، نو د $\Gamma' \Sequent !C$ يو عادي
\Log{N2i}-!!{proof}~$\delta$ شته، چې دلته د پايلې فارمول
د مخکني حکم د ښي اړخ يوازنے فارمول، او په تش ښي اړخ کښې
د دروغو فارمول دے۔ د نښه لرونکو !!{formula}s سېټ~$\Gamma'$
له څوګوني سېټ~$\Gamma$ سره مطابقت لري: يعنې د هر
!!{formula}~$!A$ دپاره په~$\Gamma'$ کښې د $x : !A$ په بڼه
د !!{element}s شمېر په~$\Gamma$ کښې د~$!A$ له تکراري شمېر،
يعنې د~$!A$ د نقلونو له شمېر، څخه کم يا ورسره برابر وي۔
\textbf{د سرچينه‌يي حد سمون OLCMP-178:} د لاندې کتنې
بې‌پرېکونه حد او د ښي اړخ د پايلې پېژندنه واضح دي؛ د کاپي
شوې کلاسيکي بڼې پر ځاے هماغه حدسي طبيعي استنتاج نوم کارېږي۔
\end{cor}

\begin{proof}
د \olref[nat][tgi]{prop:G2i-to-N2i} د ثبوت په کتنه: د
!!{proof}s په پاے کښې !!{introduction} قواعد ورزياتوو، او
!!{elimination} قواعد يوازې د !!{proof}s په پاسنۍ برخه کښې
ورزياتوو۔ نو هېڅکله !!{introduction} قاعده د
!!a{elimination} قاعدې د پاسه نۀ راولو۔
\end{proof}

خو په \olref[nat][tgi]{prop:G2i-to-N2i} کښې د \Cut{} قاعدې
ورکړل شوے اړول په \Log{N2i}-!!{proof}~$\delta$ کښې
پرېکونونه پيدا کولے شي۔ که \Log{G2i} !!{proof} په~\Cut{}
پاے ته رسېږي او نږدې فرعي !!{proof}s ئې $\pi_1$ او $\pi_2$
په ترتيب سره $\Gamma_1 \Sequent !A$ او
$!A, \Gamma_2 \Sequent \Delta_2$ ثابتوي، نو د~$\pi_1$
اړول شوے $\delta_1$ د~$\pi_2$ په اړول شوي~$\delta_2$ کښې
د $x:!A \Sequent !A$ بديهيتونو پر ځاے پېوندول کېږي۔ که
$\delta_1$ د !!a{introduction} په قاعده پاے ته رسېږي چې
اصلي !!{formula} ئې~$!A$ دے، او په~$\delta_2$ کښې بديهيت
$x:!A \Sequent !A$ د !!a{elimination} قاعدې لويه مقدمه
وي، په~$\delta$ کښې پرېکون پيدا کېږي۔

په \Log{N2i} کښې !!{proof}s په~\Log{G2i} کښې
!!{proof}s ته هم اړولے شو۔ د \olref[nat][tni]{prop:N2-to-G2}
په ثبوت کښې د \Elim{\land}، \Elim{\lif}، \Elim{\lnot}
او \Elim{\lforall} اړول په حاصل شوي \Log{G2i} !!{proof}
کښې~\Cut{} استنتاجونه راولي۔ خو که هغه \Log{N2i}-!!{proof}
چې ترې پيل کوو عادي وي، د دې مخنيوے کېدے شي۔
\textbf{د سرچينه‌يي شرط سمون OLCMP-179:} په نيمګړې جمله
کښې د عادي والي شرط پاتې و؛ د راتلونکې قضیې هماغه شرط
واضح شو۔

په \Log{N2i}-!!{proof}~$\delta$ کښې د سيکوېنټ د
موردونو لړۍ $S_1$، \dots،~$S_n$ \emph{څانګه} بولو که
$S_1$ بديهيت وي، او هر وخت چې $S_i$ د استنتاج لويه مقدمه
وي، $S_{i+1}$ ئې پايله وي۔ په بله وينا، څانګه په~$\delta$
کښې له بديهيتونو سيکوېنټونه لاندې تعقيبوي، خو د
!!{elimination} قواعدو پر کوچنيو مقدمو درېږي۔ د~$\delta$
\emph{اصلي څانګه} هغه څانګه ده چې $S_n$ ئې پاييز سيکوېنټ
وي۔ هر !!{proof}~$\delta$ لږ تر لږه يوه اصلي څانګه لري،
ځکه هره قاعده لږ تر لږه يوه لويه مقدمه لري؛ يوازې
\Intro{\land} دوه لويې مقدمې لري۔

\begin{prop}\ollabel{prop:normal-N2i-to-G2i}
که $\delta$ د~$\Gamma \Sequent !A$ عادي $\Log{N2c}$ يا
$\Log{N2i}$-!!{proof} وي، نو د $\Gamma' \Sequent \Delta$
بې‌پرېکونه $\Log{G2c}$-!!{proof}، يا په ترتيب سره
$\Log{G2i}$-!!{proof}~$\pi$، شته۔ دلته $\Gamma'$ هغه
د !!{formula}s څوګونے سېټ دے چې له $\Gamma$ څخه د نښو په
لرې کولو حاصلېږي؛ که $!A$~د~$\lfalse$ فارمول نۀ وي،
$\Delta = \{!A\}$، او که وي، $\Delta = \emptyset$ دے۔
\textbf{د سرچينې د ثبوت د حد څرګندونه OLCMP-183:} عمومي
قضيه عين ساتل شوې؛ لاندې اصلي څانګه‌يي دليل د حدسي حالت
دپاره دے۔ کلاسيکي د دروغو قاعده هم فرضونه باسي، نو د
فرضونو د نۀ ايستلو لاندې کتنه پرې نۀ تطبيقېږي۔ جلا کلاسيکي
حالت دلته نۀ دے ثابت شوے۔
\end{prop}

\begin{proof}
دا ځل استقرا د $\Gamma \Sequent !A$ په !!{proof}~$\delta$
کښې د استنتاجونو پر شمېر کوو۔ که په~$\delta$ کښې استنتاج
نۀ وي، دا بديهيت $x:!A \Sequent !A$ دے، او $\pi$ اړوند
بديهيت $!A \Sequent !A$ دے، يا که $!A \ident \lfalse$
وي، $\lfalse \Sequent \ $ دے۔

که $\delta$ په استنتاجي قاعده پاے ته رسېږي، حالتونه بېلوو:
\begin{enumerate}
  \item که $R$ د !!a{introduction} قاعده يا~\FalseInt وي،
  اړول کټ مټ د \olref[nat][tni]{prop:N2-to-G2} د ثبوت په
  شان کېږي او~\Cut ته اړتيا نۀ لري۔

  که د~$\delta$ وروستۍ قاعده~$R$ د !!a{elimination}
  قاعده وي، لاندې خبرې وګورئ:
  \begin{enumerate}
    \item د~$\delta$ له اصلي څانګې هېڅ !!{introduction}
    قاعده نۀ تېرېږي؛ کنه د اصلي څانګې په اوږدو کښې به د
    !!a{introduction} قاعدې يا~\FalseInt{} ورپسې
    !!a{elimination} قاعده راتله، او $\delta$ به عادي نۀ و۔
    \item ځکه د~$\delta$ په اصلي څانګو کښې
    !!{introduction} قاعدې نشته، يوازې يوه اصلي څانګه ده۔
    يوازې \Intro{\land} دوه لويې مقدمې لري، نو څو اصلي
    څانګې به د \Intro{\land} قاعدې له مقدمو تېرېدې۔
    \item که په اصلي څانګه کښې د \Elim{\lor} يا
    \Elim{\lexists} لويه مقدمه وي، داسې قاعده يوازې يوه ده
    او همدا وروستۍ، يعنې~$R$ ده۔ کنه په اصلي څانګه کښې به
    داسې \Elim{\lor} يا \Elim{\lexists} قاعده وه چې پايله
    ئې د !!a{elimination} قاعدې لويه مقدمه وي؛ نو ثبوت به
    عادي نۀ و، ځکه د ځاے بدلولو بدلون به تطبيقېده۔
    \item له \Elim{\lor} او \Elim{\lexists} پرته نور
    !!{elimination} قواعد فرضونه نۀ باسي، يعنې د سيکوېنټ
    کيڼ سياق نۀ کموي۔ \Elim{\lor} او \Elim{\lexists}
    يوازې د کوچنۍ مقدمې د پاسه فرضونه باسي، يعنې يوازې د
    خپلو کوچنيو مقدمو کيڼ سياقونه کموي۔ نو که د~$\delta$
    په اصلي څانګه کښې سيکوېنټ $S_i$ سياق~$\Gamma_1$ ولري،
    د هغه استنتاج پايله~$S_{i+1}$ چې $S_i$ ئې لويه مقدمه
    ده، سياق $\Gamma_1' \supseteq \Gamma_1$ لري۔
    \item له دې امله، که $x: !C \Sequent !C$ د اصلي
    څانګې تر ټولو پاسنے سيکوېنټ~$S_1$ وي، نو
    $x:!C \in \Gamma$ دے۔
  \end{enumerate}
  اوس د~$!C$ د بڼې له مخې حالتونه بېلوو۔ ځکه په اصلي
  څانګه کښې هر $S_i$ د !!a{elimination} قاعدې لويه مقدمه
  ده، دا خبره پر $x:!C \Sequent !C$ هم تطبيقېږي۔

  \item $!C \ident !D \land !E$۔ بيا $\delta$ دا بڼه لري:
\[
  \Axiom$x:!D \land !E \fCenter !D \land !E$
  \RightLabel{\Elim{\land}[1]}
  \UnaryInf$x:!D \land !E \fCenter !D$
  \Deduce$x:!D \land !E, \Gamma_1 \fCenter !A$
  \DisplayProof
  \]
  په اصلي څانګه کښې $x:!D \land !E \Sequent !D$ راځي۔
  هلته د \Intro{\lif} يا \Intro{\lnot} لويه مقدمه، يا د
  \Elim{\lor} يا \Elim{\lexists} کوچنۍ مقدمه نشته؛ نو
  د $x:!D \land !E$ د سياق !!{formula}s د اصلي څانګې په
  اوږدو کښې پاتې کېږي، که څۀ هم له کوچنيو څانګو نور فرضونه
  ورزياتېدې شي۔ له همدې امله پاييز سياق~$x:!D \land !E$
  لري۔ سرچينه وړانديز کوي چې !!{proof} $\delta$ په
  !!{proof}~$\delta_1$ داسې بدل کړو: په \Elim{\land}
  پاے ته رسېدونکے فرعي !!{proof} په $x:!D \Sequent !D$
  بدل کړو او د اصلي څانګې د هر سيکوېنټ په سياق کښې
  $x:!D \land !E$ په $x:!D$ بدل کړو، يعنې دا بدلون:
\[
  \bottomAlignProof
  \Axiom$x:!D \land !E \fCenter !D \land !E$
  \RightLabel{\Elim{\land}[1]}
  \UnaryInf$x:!D \land !E \fCenter !D$
  \Deduce$x:!D \land !E, \Gamma_1 \fCenter !A$
  \DisplayProof
  \quad\text{په دې}\quad
  \bottomAlignProof
  \Axiom$x:!D \fCenter !D$
  \Deduce$x:!D, \Gamma_1 \fCenter !A$
  \DisplayProof
  \]
  که وړانديز شوے بدلون صحيح فرعي ثبوت راکړي، د استقرا
  فرضيه پر $\delta_1$ تطبيقېږي، ځکه په $\delta$
  کښې د استنتاجونو شمېر په $\delta_1$ کښې له شمېر څخه
  په~$1$ زيات دے۔
  \textbf{د سرچينه‌يي نوم سمون OLCMP-180:} د کم شوي
  فرعي ثبوت ناتعريف پرېم لرونکے نوم د هماغه تعريف شوي
  لومړي فرعي ثبوت په نوم بدل شو۔
  \textbf{د سرچينه‌يي جوړونې د حد څرګندونه OLCMP-181:}
  د سياق ټول غړي ضرور عين نۀ پاتې کېږي؛ کوچنۍ څانګې ئې
  زياتولے شي۔ پاسنے عمومي بدلون هغه وخت بې‌نورې کتنې صحيح
  نۀ دے چې هماغه پخوانے نښه لرونکے فرض په کوچنۍ څانګه کښې
  هم استعمالېږي: د نوې مقدمې فارمول له هماغې نښې سره ټکر
  کوي، يا د اړين پخواني فرض د حذفېدو خطر شته۔ رسمي رسم
  د سرچينې وړانديز دے، د دې عمومي حالت بشپړ سم ثبوت نۀ دے
  تاييد شوے۔

  د استقرا د فرضيې له مخې، د $!D, \Gamma_1' \Sequent !A$
  يو \Log{G2i}-!!{proof} $\pi_1$ شته۔
  د \LeftR{\land} په تطبيق !!{proof}~$\pi$ اخلو:
\[
  \AxiomC{}
  \RightLabel{$\pi_1$}
  \Deduce$!D, \Gamma_1' \fCenter !A$
  \RightLabel{\LeftR{\land}}
  \UnaryInf$!D \land !E, \Gamma_1' \fCenter !A$
  \DisplayProof
  \]
  که $!A \ident \lfalse$ وي، \RightR{\Weakening}
  ورزياتوو، مثلاً:
\[
  \AxiomC{}
  \RightLabel{$\pi_1$}
  \Deduce$!D, \Gamma_1' \fCenter $
  \RightLabel{\LeftR{\land}}
  \UnaryInf$!D \land !E, \Gamma_1' \fCenter $
  \RightLabel{\RightR{\Weakening}}
  \UnaryInf$!D \land !E, \Gamma_1' \fCenter !A$
  \DisplayProof
  \]
  \item که $!C \ident !D \lor !E$ وي، ضرور د
  \Elim{\lor} لويه مقدمه ده او دا استنتاج په اصلي څانګه
  کښې وروستے استنتاج~$R$ دے۔ يعنې $\delta$ دا بڼه لري:
\[
  \Axiom$x:!D \lor !E \fCenter !D \lor !E$
  \AxiomC{}
  \RightLabel{$\delta_1$}
  \Deduce$y:!D, \Gamma_1 \fCenter!A$
  \AxiomC{}
  \RightLabel{$\delta_2$}
  \Deduce$z:!E, \Gamma_2 \fCenter !A$
  \RightLabel{\Elim{\lor}}
  \TrinaryInf$x:!D \lor !E, \Gamma_1, \Gamma_2 \fCenter !A$
  \DisplayProof
  \]
  د استقرا فرضيه پر !!{proof}s $\delta_1$ او $\delta_2$
  تطبيقېږي؛ د $!D, \Gamma_1' \Sequent !A$ او
  $!E, \Gamma_2' \Sequent !A$ دپاره \Log{G2i}-!!{proof}s
  $\pi_1$ او $\pi_2$ اخلو۔ د~$\pi$ د ترلاسه کولو دپاره
  $\LeftR{\lor}$ تطبيقوو:
\[
  \AxiomC{}
  \RightLabel{$\pi_1$}
  \Deduce$!D, \Gamma_1' \fCenter !A$
  \AxiomC{}
  \RightLabel{$\pi_2$}
  \Deduce$!E, \Gamma_2' \fCenter !A$
  \RightLabel{\LeftR{\lor}}
  \BinaryInf$!D \lor !E, \Gamma_1', \Gamma_2' \fCenter !A$
  \DisplayProof
  \]
  که $\Elim{\lor}$ فرض $y: !D$ يا $z: !E$ نۀ باسي، يعنې
  د کوچنيو مقدمو په سياق کښې دا فرضونه نۀ وي، يا که~$!A$
  د~$\lfalse$ فارمول وي، بايد ځينې \LeftR{\Weakening}
  او \RightR{\Weakening} ورزياتې کړو، مثلاً:
\[
  \AxiomC{}
  \RightLabel{$\pi_1$}
  \Deduce$\Gamma_1' \fCenter $
  \RightLabel{\LeftR{\Weakening}}
  \UnaryInf$!D, \Gamma_1' \fCenter $
  \AxiomC{}
  \RightLabel{$\pi_2$}
  \Deduce$\Gamma_2' \fCenter $
  \RightLabel{\LeftR{\Weakening}}
  \UnaryInf$!E, \Gamma_2' \fCenter $
  \RightLabel{\LeftR{\lor}}
  \BinaryInf$!D \lor !E, \Gamma_1', \Gamma_2' \fCenter $
  \RightLabel{\RightR{\Weakening}}
  \UnaryInf$!D \lor !E, \Gamma_1', \Gamma_2' \fCenter !A$
  \DisplayProof
  \]
  \textbf{د سرچينه‌يي سياق سمون OLCMP-180:} د لومړي
  حدسي سيکوېنټ ثبوت په نثر کښې هم بې‌نښو فارمولونو او
  اړوند پرېم لرونکي څوګوني سياق سره ليکل شو، لکه اصلي رسم۔

  \item $!C \ident !D \lif !E$۔ بيا $\delta$ دا بڼه لري:
\[
  \Axiom$x:!D \lif !E \fCenter !D \lif !E$
  \AxiomC{}
  \RightLabel{$\delta_1$}
  \Deduce$\Gamma_1 \fCenter!D$
  \RightLabel{\Elim{\lif}}
  \BinaryInf$x:!D \lif !E, \Gamma_1 \fCenter !E$
  \RightLabel{$\delta_2$}
  \Deduce$x:!D \lif !E, \Gamma_1, \Gamma_2 \fCenter !A$
  \DisplayProof
  \]
  پام وکړئ چې په اصلي څانګه کښې $x:!D \lif !E,
  \Gamma_1 \Sequent !E$ راځي؛ د \Intro{\lif} يا
  \Intro{\lnot} لويه مقدمه، يا د \Elim{\lor} يا
  \Elim{\lexists} کوچنۍ مقدمه نۀ لري۔ نو د
  $x:!D \lif !E, \Gamma_1$ د سياق !!{formula}s د اصلي
  څانګې په اوږدو کښې پاتې کېږي؛ نور فرضونه ورزياتېدې شي۔
  له دې امله پاييز سياق~$x:!D \lif !E, \Gamma_1$ لري۔
  سرچينه وړانديز کوي چې د !!{proof} ټوټه $\delta_2$ په
  صحيح !!{proof}~$\delta_2'$ داسې بدله شي: په
  \Elim{\lif} پاے ته رسېدونکے فرعي !!{proof} په
  $x:!E \Sequent !E$ بدل شي او د اصلي څانګې د هر سيکوېنټ
  په سياق کښې $x:!D \lif !E, \Gamma_1$ په $x:!E$
  بدل شي، يعنې:
\[
  \bottomAlignProof
  \Axiom$x:!D \lif !E, \Gamma_1 \fCenter !E$
  \RightLabel{$\delta_2$}
  \Deduce$x:!D \lif !E, \Gamma_1, \Gamma_2 \fCenter !A$
  \DisplayProof
  \quad\text{په دې}\quad
  \bottomAlignProof
  \Axiom$x:!E \fCenter !E$
  \RightLabel{$\delta_2'$}
  \Deduce$x:!E, \Gamma_2 \fCenter !A$
  \DisplayProof
  \]
  که وړانديز شوے بدلون صحيح فرعي ثبوت راکړي، د استقرا
  فرضيه پر $\delta_1$ او $\delta_2'$ تطبيقېږي،
  ځکه په $\delta$ کښې د استنتاجونو شمېر په $\delta_1$ او
  ~$\delta_2$ کښې د استنتاجونو د شمېرونو مجموعه ده، او
  ~$1$ د اصلي څانګې په پيل کښې د \Elim{\lif} استنتاج
  دپاره ورسره زياتېږي۔ که همدغه استنتاج په $\delta$ کښې
  وروستے استنتاج~$R$ هم وي، نو $\delta_2$ د $0$
  استنتاجونو ټوټه ده۔
  \textbf{د سرچينه‌يي شمېر سمون OLCMP-182:} په وروستي
  استنتاج کښې بې‌استنتاجه پاتې ټوټه دويمه ده؛ کوچنۍ مقدمه
  ضرور بې‌استنتاجه نۀ وي۔ د پخوانيو فرضونو د حذفولو
  OLCMP-181 حد په دې وړانديز هم تطبيقېږي۔

  د استقرا د فرضيې له مخې، د $\Gamma_1' \Sequent !D$
  او $!E, \Gamma_2' \Sequent !A$ دپاره په ترتيب سره
  \Log{G2i}-!!{proof}s $\pi_1$ او $\pi_2$ شته۔
  د \LeftR{\lif} په تطبيق !!{proof}~$\pi$ اخلو:
\[
  \AxiomC{}
  \RightLabel{$\pi_1$}
  \Deduce$\Gamma_1' \fCenter !D$
  \AxiomC{}
  \RightLabel{$\pi_2$}
  \Deduce$!E, \Gamma_2' \fCenter !A$
  \RightLabel{\LeftR{\lif}}
  \BinaryInf$!D \lif !E, \Gamma_1', \Gamma_2' \fCenter !A$
  \DisplayProof
  \]
  که $!D$ د دروغو فارمول وي، او/يا $!A \ident \lfalse$
  وي، يو يا دوه \RightR{\Weakening} ورزياتوو، مثلاً:
\[
  \AxiomC{}
  \RightLabel{$\pi_1$}
  \Deduce$\Gamma_1' \fCenter $
  \RightLabel{\RightR{\Weakening}}
  \UnaryInf$\Gamma_1' \fCenter !D$
  \AxiomC{}
  \RightLabel{$\pi_2$}
  \Deduce$!E, \Gamma_2' \fCenter $
  \RightLabel{\RightR{\Weakening}}
  \UnaryInf$!E, \Gamma_2' \fCenter !A$
  \RightLabel{\LeftR{\lif}}
  \BinaryInf$!D \lif !E, \Gamma_1', \Gamma_2' \fCenter !A$
  \DisplayProof
  \]
  \textbf{د سرچينه‌يي سياق سمون OLCMP-182:} په نثر کښې
  د لومړي ثبوت سياق پرېم لري او د دويم ثبوت سياق د دويمې
  ټوټې دے، لکه په دواړو اصلي رسمونو کښې۔

\end{enumerate}
پاتې حالتونه په ورته طريقه څېړل کېږي او د مشقونو په توګه
پرېښودل شوي دي۔
\end{proof}

\begin{cor}
په عادي \Log{N1i} يا \Log{N2i}-!!{proof} کښې هر راغلے
!!{formula} د پاييز !!{formula} يا د پاتې پرانيستو فرضونو
فرعي‌!!{formula}، يا د پاييز سيکوېنټ اړوند فرعي فارمول دے؛
په لومړي ترتيب کښې اړوند کمي تعويضي موردونه هم مراد دي۔
\textbf{د سرچينه‌يي حد سمون OLCMP-184، د OLCMP-169 تکراري مورد:}
د طبيعي استنتاج يوازې وروستے فارمول بس نۀ دے؛ پاتې
پرانيستي فرضونه او اړوند کمي موردونه هم په دې حد کښې راځي۔
\end{cor}

\begin{proof}
د \olref{prop:normal-N2i-to-G2i} له مخې، هر عادي
\Log{N2i}-!!{proof} په \Log{G2i}-!!{proof} اړولے شو۔ د
ثبوت کتنه ښيي چې په \Log{N2i}-!!{proof} کښې هر راغلے
!!{formula}، د دروغو د پايلې له استثنا سره، په
\Log{G2i}-!!{proof} کښې هم راځي؛ د دروغو پايله دلته په
تش ښي اړخ ښودل کېږي۔ \Log{G2i} د \Cut{} قاعده نۀ لري،
نو د اړوند فرعي‌!!{formula} خاصيت لري۔ د اصلي جوړونې
پورته څرګند شوي پاتې حالتونه د بشپړې کتنې ځاے نۀ نيسي۔
\end{proof}

\end{document}
